% Syntra LaTeX Template
% For Arabic support, use XeLaTeX with:
%   \usepackage{fontspec}
%   \usepackage{polyglossia}
%   \setdefaultlanguage{english}
%   \setotherlanguage{arabic}
%   \newfontfamily\arabicfont[Script=Arabic]{Amiri}

\documentclass[12pt,a4paper]{article}

% Font and encoding
\usepackage[utf8]{inputenc}
\usepackage[T1]{fontenc}

% Math packages
\usepackage{amsmath,amssymb}

% Table support
\usepackage{longtable,booktabs}

% Graphics and links
\usepackage{graphicx}
\usepackage{hyperref}
\usepackage{geometry}

% Page geometry
\geometry{
    left=2.5cm,
    right=2.5cm,
    top=3cm,
    bottom=3cm
}

% Fix pandoc tightlist issue
\providecommand{\tightlist}{\setlength{\itemsep}{0pt}\setlength{\parskip}{0pt}}

% Hyperref settings
\hypersetup{
    colorlinks=true,
    linkcolor=blue,
    filecolor=magenta,
    urlcolor=cyan,
    pdftitle={Syntra Research Documentation},
    pdfauthor={Syntra-Env},
    pdfpagemode=FullScreen
}

\begin{document}

\section{Search-Write Juxtaposition}\label{search-write-juxtaposition}

\subsection{Overall Interaction Model}\label{overall-interaction-model}

Scholar should be organized around two primary display modes---Reading
Mode and Writing Mode---each corresponding to a distinct attentional
regime and explicitly designed to regulate cognitive load.

\subsection{Reading Mode}\label{reading-mode}

In Reading Mode, the Quranic text is the exclusive foreground. The
writing pane is fully suppressible and hidden by default, enforcing a
strong minimal baseline in which no analytical information is surfaced
unless explicitly requested. Writing remains one deliberate action away
through a persistent affordance (for example, a toggle or keyboard
shortcut), allowing interpretive intervention without navigational
transitions, mode confusion, or loss of textual context.

\subsection{Writing Mode}\label{writing-mode}

In Writing Mode, the interface resolves into two synchronized panes: the
Quranic text on one side and a persistent note-writing surface on the
other. Writing is treated as a first-order analytical activity rather
than an auxiliary feature. The writing surface remains available across
reading, navigation, and search operations, enabling the capture of
observations, hypotheses, interpretive claims, and methodological
reflections at the moment they arise, while maintaining a stable spatial
separation between text, interpretation, and evidence.

\subsection{Multi-Resolution Annotation
Model}\label{multi-resolution-annotation-model}

Notes are not merely attached to passages; they are attached to
addressable analytical objects at multiple resolutions. At minimum, the
system supports attachment to an ayah, to a token instance within an
ayah, to a word-form or morphological realization, to a root or
lexeme-level abstraction, and to syntactic objects at both the type
level (a general construction or pattern) and the instance level (a
particular occurrence of that construction in context). The interface
explicitly distinguishes between type-level annotations (general claims)
and instance-level annotations (context-bound claims), and never merges
them implicitly. Traversal between levels is always user-initiated and
scope-labeled.

\subsection{Cognitive Load and Scoping
Defaults}\label{cognitive-load-and-scoping-defaults}

Cognitive overload is controlled through strict scoping defaults. When a
user reads an ayah, the notes view surfaces only notes directly attached
to that ayah. Related notes attached to contained objects (token
instances, syntactic instances) are hidden by default and may be
revealed only through explicit, reversible scope expansion. Expanded
notes are grouped by attachment level and visually labeled so the user
can always determine why a note is visible.

\subsection{Lexical Generalization Without Loss of
Specificity}\label{lexical-generalization-without-loss-of-specificity}

When a user inspects a word at the root or lexeme layer, the system
preserves specificity by presenting word-form annotations as
subordinate, form-scoped notes within the broader root-level inspection.
Generalization never collapses detail; instead, higher-level views
summarize lower-level notes and allow drill-down into concrete forms and
occurrences without forcing their simultaneous display.

\subsection{Syntactic Patterns and Evidence
Separation}\label{syntactic-patterns-and-evidence-separation}

For syntactic elements, the UI enforces a separation between
interpretation and evidence. Opening the notes for a syntactic type
(pattern) presents interpretive notes in one pane, while the adjacent
pane displays a bounded, non-exhaustive list of concrete instances
satisfying that pattern. Notes may be attached both to the pattern
itself and to individual instances, but these attachments remain
visually and structurally distinct. Instance lists are summarized,
capped, or aggregated to prevent unbounded expansion.

\subsection{Navigability Through Explicit Scope
Control}\label{navigability-through-explicit-scope-control}

Navigability is preserved through scoping rather than proliferation. By
default, the user sees only direct notes for the currently selected
object. Explicit affordances allow scope expansion upward to parent
abstractions (instance → type → schema) or downward to contained
refinements (root → word-forms → token instances), with each expansion
clearly labeled and independently reversible.

\subsection{Search as an Instrument of
Inquiry}\label{search-as-an-instrument-of-inquiry}

Search operates over a formally specified set of searchable elements,
including lemma, root, morphological structure, syntactic annotation,
phonological markers, stylistic features, and higher-order pattern
signatures. Each validated search yields a stable query object that can
itself be annotated and later revisited. Search results are treated as
evidence sets rather than content streams, reinforcing
Scholar\textquotesingle s role as an instrument for inquiry rather than
a retrieval interface. Throughout, Reading Mode preserves a
distraction-free baseline, and Writing Mode exposes analytical power
only through deliberate, scoped interaction.

\end{document}
